% Generated draft: requires mathematical and source review.
\documentclass[11pt]{article}
\usepackage[margin=25mm]{geometry}
\usepackage{fontspec}
\setmainfont{Latin Modern Roman}
\usepackage{amsmath,amssymb,mathtools,unicode-math}
\setmathfont{Latin Modern Math}
\usepackage{enumitem,needspace,xurl,hyperref}
\setlength{\parindent}{0pt}
\setlength{\parskip}{5pt}
\setlength{\emergencystretch}{3em}
\newcommand{\N}{\mathbb{N}}
\newcommand{\Z}{\mathbb{Z}}
\newcommand{\Q}{\mathbb{Q}}
\newcommand{\R}{\mathbb{R}}
\newcommand{\C}{\mathbb{C}}
\newcommand{\F}{\mathbb{F}}
\newcommand{\A}{\mathbb{A}}
\newcommand{\PP}{\mathbb{P}}
\newcommand{\E}{\mathbb{E}}
\newcommand{\eps}{\varepsilon}
\newcommand{\dd}{\,\mathrm{d}}
\begin{document}
\section*{UnsolvedMath ID: 30000232}
\textbf{Source problem code:} catalogue entry
\subsection{Definitions and mathematical statement}
Given an initial rod configuration $\Gamma\subset\R^3$, the source asks whether geometric gradient flows can find critical points of a Kirchhoff elastic energy, called elasticae in the passage. A critical point is stationary under all variations admissible for the chosen rod model and constraints; it need not be a global minimum, and finding one is different from enumerating every critical point.

A centerline is a regular curve $\gamma$ with arclength parameter $s$, unit tangent $\tau=\partial_s\gamma$ and curvature $\kappa=|\partial_s\tau|$. A full Kirchhoff rod can also require a material orthonormal frame $(d_1,d_2,d_3)$ with $d_3=\tau$. Writing $\partial_sd_i=\Omega\times d_i$ defines bending/twist strains; the twist component is $\Omega\cdot d_3$ and need not equal the Frenet torsion. Thus an unframed point set $\Gamma$ alone does not determine a general twist energy.

For example, an isotropic framed model can have energy
\[
 E(q)=\frac12\int\bigl(a\kappa^2+c(\Omega\cdot d_3)^2\bigr)\,ds,
 \qquad a,c>0,
\]
where $q$ includes the curve and frame. This is an explicit illustrative model, not the report's recovered complete functional: intrinsic strain, anisotropy, forces, endpoint conditions, length and topology constraints are not specified in the extracted question.

After fixing a configuration space $\mathcal M$ and an appropriate metric, the desired constrained gradient evolution has the form
\[
 \partial_tq=-\operatorname{grad}_{\mathcal M}E(q),\qquad q(0)=q_0.
\]
For a sufficiently regular flow its intended descent identity is $\frac{d}{dt}E(q(t))=-\|\operatorname{grad}_{\mathcal M}E(q(t))\|^2$. The substantive request is to construct a well-defined flow that remains admissible for long time and converges, in an appropriate sense and possibly modulo rigid motion/reparametrization, to an equilibrium. Energy decrease alone is not an existence or convergence theorem. Whether one requires full-sequence convergence, subsequential convergence or a numerical algorithm with guarantees is not stated.

The supplied literature triage checked 21 August 2026 reports long-time length-preserving flows and convergence to equilibrium for isotropic Kirchhoff models, including embedded-flow results. It distinguishes those theorems from an unavailable general algorithm for arbitrary three-dimensional configurations and energies. Their exact admissibility, framing, boundary and self-intersection hypotheses are not reproduced; no extension beyond their reported scope is asserted.
\subsection{Short English statement}
Can a geometric gradient flow starting from a spatial rod approach an equilibrium of a Kirchhoff elastic energy? The energy, constraints, frame and convergence notion must be fixed. The supplied assessment reports long-time and convergence results for isotropic length-preserving models, without a general arbitrary-energy algorithm.
\subsection{Sources}
\begin{itemize}
\item Local UnsolvedMath registry, permanent ID 30000232, exact source code catalogue entry; \emph{Geometric Gradient Flows for Elastic Rods}. Proposed year 2005, status partially\_solved; corrected\_verified extraction removes only report lead-ins and restores spatial notation/grammar.
\item Supplied report citation: \emph{Nonlinear Evolution Problems (2006)}, Oberwolfach Reports, DOI 10.4171/owr/2005/25, pp. 1371--1430, \url{https://doi.org/10.4171/owr/2005/25}. Title year differs from DOI/proposed year.
\item Supplied literature triage checked 2026-08-21: ``original Oberwolfach Report 25/2005'', \url{https://doi.org/10.4171/owr/2005/25}; ``Lin--Schwetlick, geometric flow of Kirchhoff rods'', \url{https://doi.org/10.1137/S0036139903431713}; ``Lin--Schwetlick, rods without self-intersections'', \url{https://doi.org/10.1007/s10910-008-9383-6}. This dated OWR-derived assessment requests independent verification of time-sensitive status and cited claims.
\end{itemize}
\end{document}
